\section{Conditional contraction and identifiability}
\label{sec:contraction}

Well-posedness for a fixed $\sigobs$ does not determine what happens
as information increases.  A contraction rate is a property of a
sequence of statistical experiments, not of the prior alone.  This
section records a standard criterion in a form that makes the missing
inputs explicit.

\subsection{A test-and-prior-mass criterion}

For each $n$, let
$\{P_u^{(n)}:u\in\HH\}$ be a dominated experiment, let
$Y^{(n)}\sim P_{\udag}^{(n)}$, and let $\Pi_n(\,\cdot\mid Y^{(n)})$
be the posterior obtained from the fixed prior $\muz$.  Let $d_n$ be
a semimetric.  Whenever the relevant Radon--Nikodym derivatives exist,
write
\begin{align}
  K_n(u)
    &=\KL\!\left(P_{\udag}^{(n)}\,\|\,P_u^{(n)}\right),\\
  V_n(u)
    &=\operatorname{Var}_{P_{\udag}^{(n)}}\!
      \left[
        \log\frac{\mathrm dP_{\udag}^{(n)}}{\mathrm dP_u^{(n)}}
      \right].
\end{align}

\begin{theorem}[Conditional contraction criterion]
\label{thm:contraction}
Let $\epsilon_n\downarrow0$ with $n\epsilon_n^2\to\infty$.  Suppose
there are a constant $c>0$, measurable sieves
$\mathcal F_n\subset\HH$, measurable sets
\begin{equation}
  \mathcal K_n
  \subset
  \{u:K_n(u)\leq n\epsilon_n^2,\;
       V_n(u)\leq n\epsilon_n^2\},
  \label{eq:kl-neighborhood}
\end{equation}
tests $\varphi_n$, and a constant $M<\infty$ such that
\begin{align}
  \muz(\mathcal K_n)
    &\geq e^{-c n\epsilon_n^2}, \label{eq:contraction-mass}\\
  \muz(\mathcal F_n^c)
    &\leq e^{-(c+4)n\epsilon_n^2}, \label{eq:contraction-sieve}\\
  P_{\udag}^{(n)}\varphi_n
    &\longrightarrow0, \label{eq:contraction-type1}\\
  \sup_{\substack{u\in\mathcal F_n\\
                   d_n(u,\udag)>M\epsilon_n}}
  P_u^{(n)}(1-\varphi_n)
    &\leq e^{-(c+4)n\epsilon_n^2}. \label{eq:contraction-type2}
\end{align}
Then
\begin{equation}
  \Pi_n\!\left(
    u:d_n(u,\udag)>M\epsilon_n
    \,\middle|\,Y^{(n)}
  \right)
  \longrightarrow0
  \quad\text{in }P_{\udag}^{(n)}\text{-probability}.
  \label{eq:contraction-conclusion}
\end{equation}
\end{theorem}

This theorem is a direct test-and-denominator argument; see
\Cref{sec:appendix-contraction}.  Entropy conditions often appear in
applications because they are a way to construct
\eqref{eq:contraction-type1}--\eqref{eq:contraction-type2}.  They are
not consequences of the drift assumptions.

\subsection{What Girsanov does and does not transfer}

Suppose \Cref{ass:girsanov} holds, so that
$\muz=\rho\,\mu_{\mathrm{ref}}$.

\begin{corollary}[Local mass transfer]
\label{cor:local-mass-transfer}
In the setting of \Cref{thm:contraction}, suppose
\begin{align}
  \mu_{\mathrm{ref}}(\mathcal K_n)
    &\geq e^{-c_0n\epsilon_n^2},\\
  \operatorname*{ess\,inf}_{u\in\mathcal K_n}
    \rho(u)
    &\geq e^{-a n\epsilon_n^2}
\end{align}
for constants $c_0,a\geq0$.  Then
\begin{equation}
  \muz(\mathcal K_n)
  \geq e^{-(c_0+a)n\epsilon_n^2}.
  \label{eq:mass-transfer}
\end{equation}
If, in addition, the sieve and testing conditions of
\Cref{thm:contraction} hold with $c=c_0+a$, then
\eqref{eq:contraction-conclusion} follows.
\end{corollary}

\begin{proof}
Integrating the local lower bound gives
\[
  \muz(\mathcal K_n)
  =\int_{\mathcal K_n}\rho\,\mathrm d\mu_{\mathrm{ref}}
  \geq e^{-a n\epsilon_n^2}
       \mu_{\mathrm{ref}}(\mathcal K_n),
\]
which is \eqref{eq:mass-transfer}.  Apply
\Cref{thm:contraction}.
\end{proof}

The hypothesis on the essential infimum is genuinely local.
The global estimate $\rho\in L^2(\mu_{\mathrm{ref}})$ from
\eqref{eq:terminal-density} gives an upper-integrability bound; it
does not prevent $\rho$ from becoming arbitrarily small on the
shrinking sets $\mathcal K_n$.  It therefore cannot be used by itself
to claim that a Gaussian contraction rate transfers to the learned
prior.

Even for Gaussian priors, there is no model-independent exponent.  In
the sequence model analysed by \citet{knapik2011}, the rate depends on
the regularity of the prior, the regularity of the truth, the
ill-posedness of the forward operator, and any $n$-dependent prior
scale.  A rate for \eqref{eq:prior-sde} can only be stated after those
quantities and the local density bound in
\Cref{cor:local-mass-transfer} have been established.

\subsection{The identifiable object}

For the observation model \eqref{eq:forward}, define
\[
  d_{\G}(u,v)
  =\norm{\Gamobs^{-1/2}\G(u-v)}_{\R^k}.
\]
This is generally a semimetric, not a norm.

\begin{proposition}[Observational equivalence]
\label{thm:ident}
For any $u,v\in\HH$ and any $\sigobs>0$, the data distributions in
\eqref{eq:forward} are equal if and only if
\begin{equation}
  \G u=\G v.
  \label{eq:observational-equivalence}
\end{equation}
If $S=\operatorname{supp}\muz$, the parameter is identifiable on $S$
if and only if $\G|_S$ is injective, equivalently
\begin{equation}
  (S-S)\cap\ker\G=\{0\}.
  \label{eq:support-identifiability}
\end{equation}
\end{proposition}

\begin{proof}
The two Gaussian distributions have the same positive-definite
covariance $\sigobs^2\Gamobs$, so they are equal exactly when their
means agree.  The second statement is the definition of injectivity
restricted to $S$.
\end{proof}

Under \Cref{ass:girsanov}, $\muz\sim\mu_{\mathrm{ref}}$ and hence
$S=\HH$.  If $\HH$ is infinite-dimensional and $k<\infty$, then
$\ker\G\neq\{0\}$.  Full $\HH$-norm recovery is therefore impossible
from this observation model without further structure.
One may instead prove contraction in $d_{\G}$, restrict the support to
a class on which $\G$ is injective, or assume a quantitative inverse
stability estimate linking $\norm{u-v}_{\HH}$ to $d_{\G}(u,v)$ on the
relevant sieve.  None of these properties follows from
\Cref{ass:drift}.
